% =============================================================================
% NoteCognition — University Project Report (Blackbook)
% =============================================================================
\documentclass[12pt, a4paper, oneside]{report}

% ---------------------------------------------------------------------------
% PACKAGES
% ---------------------------------------------------------------------------
\usepackage[utf8]{inputenc}
\usepackage[T1]{fontenc}
\usepackage{times}                         % Times New Roman font
\usepackage[margin=1in]{geometry}          % 1-inch margins
\usepackage{setspace}                      % Line spacing
\usepackage{graphicx}                      % Images
\usepackage{float}                         % Float placement [H]
\usepackage{caption}                       % Caption customization
\usepackage{subcaption}                    % Sub-figures
\usepackage{booktabs}                      % Professional tables
\usepackage{longtable}                     % Multi-page tables
\usepackage{tabularx}                      % Flexible-width tables
\usepackage{array}                         % Extended column types
\usepackage{listings}                      % Code listings
\usepackage{xcolor}                        % Colors
\usepackage{hyperref}                      % Clickable links
\usepackage{url}                           % URL formatting
\usepackage{amsmath}                       % Math
\usepackage{enumitem}                      % List customization
\usepackage{titlesec}                      % Section title formatting
\usepackage{fancyhdr}                      % Headers and footers
\usepackage{tocloft}                       % TOC customization
\usepackage{pdfpages}                      % Include external PDFs
\usepackage{tikz}                          % Diagrams
\usetikzlibrary{shapes.geometric, arrows.meta, positioning, calc}

% ---------------------------------------------------------------------------
% DOCUMENT SETTINGS
% ---------------------------------------------------------------------------
\onehalfspacing                            % 1.5 line spacing

% Hyperlink styling
\hypersetup{
    colorlinks=true,
    linkcolor=black,
    citecolor=black,
    urlcolor=blue,
    bookmarksopen=true,
}

% Code listing style
\definecolor{codebg}{RGB}{245,245,245}
\definecolor{codegreen}{RGB}{0,128,0}
\definecolor{codegray}{RGB}{128,128,128}
\definecolor{codeblue}{RGB}{0,0,180}
\definecolor{codepurple}{RGB}{128,0,128}

\lstdefinestyle{codestyle}{
    backgroundcolor=\color{codebg},
    basicstyle=\ttfamily\small,
    breaklines=true,
    captionpos=b,
    commentstyle=\color{codegreen},
    keywordstyle=\color{codeblue}\bfseries,
    stringstyle=\color{codepurple},
    numberstyle=\tiny\color{codegray},
    numbers=left,
    numbersep=8pt,
    frame=single,
    framerule=0.5pt,
    rulecolor=\color{codegray},
    showstringspaces=false,
    tabsize=2,
    xleftmargin=1.5em,
    framexleftmargin=1.5em,
}
\lstset{style=codestyle}

% Page style
\pagestyle{fancy}
\fancyhf{}
\fancyhead[L]{\leftmark}
\fancyhead[R]{\thepage}
\renewcommand{\headrulewidth}{0.4pt}

% Chapter title formatting
\titleformat{\chapter}[display]
  {\normalfont\Large\bfseries\centering}
  {\chaptertitlename\ \thechapter}{12pt}{\LARGE}
\titlespacing*{\chapter}{0pt}{-20pt}{20pt}

% ---------------------------------------------------------------------------
% METADATA — EDIT THESE VALUES
% ---------------------------------------------------------------------------
\newcommand{\projecttitle}{NoteCognition}
\newcommand{\projectsubtitle}{A Local-First, Cloud-Synced Markdown Note-Taking Platform with Host-Your-Own-Data Architecture}
\newcommand{\collegename}{Your College Name}               % <-- EDIT
\newcommand{\universityname}{Your University Name}         % <-- EDIT
\newcommand{\departmentname}{Department of Computer Engineering} % <-- EDIT
\newcommand{\academicyear}{2025--2026}                     % <-- EDIT
\newcommand{\guidename}{Prof. Guide Name}                  % <-- EDIT
\newcommand{\guidetitle}{Assistant Professor}               % <-- EDIT

% Student details (add/remove as needed)
\newcommand{\studentA}{Teja Lonkar}                        % <-- EDIT
\newcommand{\rollA}{Roll Number}                           % <-- EDIT
\newcommand{\studentB}{Student B Name}                     % <-- EDIT
\newcommand{\rollB}{Roll Number}                           % <-- EDIT
\newcommand{\studentC}{Student C Name}                     % <-- EDIT
\newcommand{\rollC}{Roll Number}                           % <-- EDIT
\newcommand{\studentD}{Student D Name}                     % <-- EDIT
\newcommand{\rollD}{Roll Number}                           % <-- EDIT

% ===========================================================================
\begin{document}

% ---------------------------------------------------------------------------
% FRONT MATTER (Roman page numbers)
% ---------------------------------------------------------------------------
\pagenumbering{roman}

% ===========================
% TITLE PAGE
% ===========================
\begin{titlepage}
    \centering
    \vspace*{0.5cm}

    % University / College Logo
    % \includegraphics[width=3cm]{images/college_logo.png}\\[0.5cm]

    {\large \textbf{\universityname}}\\[0.3cm]
    {\large \collegename}\\[0.3cm]
    {\large \departmentname}\\[1.5cm]

    \rule{\linewidth}{0.5pt}\\[0.4cm]
    {\Huge \textbf{\projecttitle}}\\[0.3cm]
    {\large \projectsubtitle}\\[0.2cm]
    \rule{\linewidth}{0.5pt}\\[1.5cm]

    {\large \textbf{A Project Report}}\\[0.3cm]
    {\large Submitted in partial fulfillment of the requirements for the degree of}\\[0.3cm]
    {\large \textbf{Bachelor of Engineering}}\\[0.3cm]
    {\large in}\\[0.3cm]
    {\large \textbf{Computer Engineering}}\\[1.5cm]

    {\large \textbf{Submitted by:}}\\[0.4cm]
    \begin{tabular}{ll}
        \studentA & \rollA \\
        \studentB & \rollB \\
        \studentC & \rollC \\
        \studentD & \rollD \\
    \end{tabular}\\[1.5cm]

    {\large \textbf{Under the Guidance of}}\\[0.3cm]
    {\large \guidename}\\[0.1cm]
    {\large \guidetitle}\\[1.5cm]

    {\large Academic Year: \academicyear}

    \vfill
\end{titlepage}

% ===========================
% CERTIFICATE
% ===========================
\chapter*{Certificate}
\addcontentsline{toc}{chapter}{Certificate}

This is to certify that the project titled \textbf{``\projecttitle''} has been successfully completed by the following students of \textbf{\departmentname}, \textbf{\collegename}, in partial fulfillment of the requirements for the degree of \textbf{Bachelor of Engineering} in \textbf{Computer Engineering} during the academic year \textbf{\academicyear}.

\vspace{1cm}
\begin{tabular}{ll}
    \textbf{Sr. No.} & \textbf{Student Name} \\
    \midrule
    1. & \studentA \\
    2. & \studentB \\
    3. & \studentC \\
    4. & \studentD \\
\end{tabular}

\vspace{2cm}
\begin{minipage}[t]{0.45\textwidth}
    \centering
    \rule{6cm}{0.4pt}\\
    \textbf{\guidename}\\
    Project Guide
\end{minipage}
\hfill
\begin{minipage}[t]{0.45\textwidth}
    \centering
    \rule{6cm}{0.4pt}\\
    \textbf{Head of Department}\\
    \departmentname
\end{minipage}

\vspace{2cm}
\begin{center}
    \rule{6cm}{0.4pt}\\
    \textbf{Principal}\\
    \collegename
\end{center}

% ===========================
% ACKNOWLEDGEMENT
% ===========================
\chapter*{Acknowledgement}
\addcontentsline{toc}{chapter}{Acknowledgement}

We would like to express our sincere gratitude to our project guide, \textbf{\guidename}, for providing invaluable guidance and support throughout the development of this project. Their technical expertise and constructive feedback were instrumental in shaping the direction of this work.

We are also grateful to the \textbf{Head of the \departmentname} and the \textbf{Principal of \collegename} for providing the necessary infrastructure and resources.

We extend our thanks to all faculty members and peers who contributed their time and knowledge during the course of this project.

Finally, we thank our families for their constant encouragement and support.

\vspace{1.5cm}
\begin{flushright}
    \studentA\\
    \studentB\\
    \studentC\\
    \studentD\\
\end{flushright}

% ===========================
% ABSTRACT
% ===========================
\chapter*{Abstract}
\addcontentsline{toc}{chapter}{Abstract}

Note-taking has evolved from pen-and-paper journals into sophisticated digital platforms that serve as the backbone of personal knowledge management for millions of students, researchers, and professionals worldwide. Modern users demand applications that are fast, available offline, seamlessly synchronized across devices, and---increasingly---respectful of their data privacy. Unfortunately, the overwhelming majority of popular note-taking solutions store user data exclusively on proprietary cloud servers controlled by the service provider. Users have no visibility into how their notes are stored, processed, or potentially monetized. The few alternatives that do prioritize local storage typically lack real-time synchronization, forcing users to choose between convenience and control. This fundamental tension between cloud-powered functionality and personal data sovereignty is the core problem that this project addresses.

NoteCognition is a full-stack, production-grade Markdown note-taking platform designed from the ground up around a local-first philosophy. In a local-first application, the user's device is treated as the primary source of truth rather than a remote server. Every interaction---whether creating a new note, editing text, reorganizing folders, or deleting content---is executed immediately against the browser's local IndexedDB database through the Dexie.js library. The result is an editing experience with zero perceptible latency: the application never waits for a server response before reflecting a user's action on screen. This architecture also means the application remains fully functional even when the user has no internet connection whatsoever. Notes can be written and organized entirely offline, and all pending changes are automatically queued and pushed to the cloud once connectivity is restored.

What sets NoteCognition apart from other local-first tools is its unique Host-Your-Own-Data model. Rather than routing user data through a centralized backend operated by the development team, NoteCognition empowers each user to deploy their own completely independent, private backend infrastructure on their personal Amazon Web Services account. This deployment is achieved through a single AWS CloudFormation template file that the user downloads directly from within the application. By uploading this template to the AWS CloudFormation Console, the user provisions---with a single click and without writing any code---a complete serverless backend stack comprising Amazon Cognito for user authentication, Amazon API Gateway for REST and WebSocket API endpoints, AWS Lambda functions for stateless business logic, an Amazon DynamoDB table for persistent storage using a highly optimized single-table design, and an Amazon S3 bucket for note attachments such as images. At no point during this process does the user share their AWS credentials or access keys with NoteCognition or any third party. The user's data lives entirely within their own AWS account, under their own billing, their own region of choice, and their own administrative control.

Once the backend is deployed, NoteCognition connects to it and initiates a bidirectional synchronization loop. Outbound synchronization is handled through incremental pushes: whenever a local database record changes, the modification is added to a persistent sync queue and transmitted to the server via the REST API. Inbound synchronization uses a timestamp-based pull mechanism that retrieves all remote changes newer than the client's last known sync point. In cases where the same note has been modified on multiple devices before a sync cycle completes, conflicts are resolved deterministically using a last-write-wins strategy that compares timestamps at millisecond precision. In parallel with the REST-based sync channel, a dedicated WebSocket connection provides real-time, character-by-character propagation of edits. When a user types on one device, the change is broadcast through the WebSocket API Gateway to all other active sessions belonging to that user within milliseconds, enabling a fluid cross-device editing experience.

The frontend of the application is built with React and bundled with Vite. It features a rich split-pane Markdown editor with a live-rendered preview, a hierarchical folder tree with drag-and-drop support, inline image attachments via drag-and-drop or clipboard paste, and a polished dark-mode glassmorphism user interface with responsive layouts, smooth animations, and modern typography. The entire system is designed to be lightweight, cost-efficient, and entirely operable within the AWS Free Tier, making it accessible to individual users and students without incurring significant cloud expenses. This report presents the complete design, development, testing, and deployment of NoteCognition in detail.

% ===========================
% TABLE OF CONTENTS
% ===========================
\tableofcontents
\listoffigures
\addcontentsline{toc}{chapter}{List of Figures}
\listoftables
\addcontentsline{toc}{chapter}{List of Tables}

% ---------------------------------------------------------------------------
% MAIN MATTER (Arabic page numbers)
% ---------------------------------------------------------------------------
\cleardoublepage
\pagenumbering{arabic}

% ===========================
% CHAPTER 1: INTRODUCTION
% ===========================
\chapter{Introduction}

\section{Background}

The way individuals capture, organize, and retrieve information has undergone a dramatic transformation over the past two decades. Traditional pen-and-paper note-taking has steadily given way to digital platforms that promise searchability, portability, and cross-device access. Applications such as Evernote, Google Keep, Notion, and Microsoft OneNote have become deeply embedded in the daily workflows of millions of students, educators, researchers, and professionals. These platforms have demonstrated that digital note-taking, when done well, can significantly improve productivity, knowledge retention, and collaborative work.

However, the rapid adoption of cloud-based note-taking has introduced a set of challenges that the industry has largely failed to address. The most pressing of these is the issue of data ownership and privacy. Nearly all mainstream note-taking platforms operate on a centralized cloud model in which user data is stored exclusively on servers owned and operated by the service provider. Users have little to no visibility into where their data physically resides, how it is encrypted at rest, whether it is scanned for advertising or analytics purposes, or what would happen to their notes if the provider discontinues the service. High-profile data breaches at companies like Evernote in 2013, which exposed fifty million user accounts, have underscored the real-world risks of entrusting personal knowledge bases to third-party servers.

At the same time, the few existing solutions that do prioritize local data storage, such as Obsidian, suffer from a different limitation: they lack built-in cloud synchronization capabilities. Users who wish to access their notes across multiple devices must configure third-party sync tools like Syncthing, iCloud, or Dropbox, adding complexity and introducing potential points of failure. Real-time collaborative or cross-device editing is virtually absent in locally stored note platforms.

This project was conceived to bridge this gap. The goal was to build a note-taking application that combines the performance and privacy advantages of local-first storage with the convenience of real-time cloud synchronization, while giving each individual user full administrative control over their own backend infrastructure. The result is NoteCognition, a Markdown-based note-taking platform built with React, Vite, and IndexedDB on the frontend and a fully serverless AWS backend that users can deploy on their own personal AWS accounts with a single click.

The technological landscape of 2025 makes this approach newly viable. Browser-based storage APIs such as IndexedDB have matured to the point where they can reliably serve as primary databases for complex web applications. The serverless paradigm on AWS, specifically the combination of Lambda, DynamoDB, API Gateway, Cognito, and CloudFormation, allows a complete production-grade backend to be defined as a single declarative template and deployed without provisioning or managing any servers. WebSocket support in API Gateway enables real-time bidirectional communication at negligible cost. These converging technologies form the foundation upon which NoteCognition is built.

\section{Objectives}

The primary objective of this project is to design and develop a full-stack, production-grade Markdown note-taking platform that operates on a local-first architecture while providing seamless real-time synchronization through a user-owned serverless cloud backend. The application must store all note content and folder structures directly in the browser's IndexedDB database so that every user interaction, whether creating a note, editing text, reorganizing folders, or deleting content, is executed instantly against the local store without blocking on any network request. This ensures that the application delivers a zero-lag editing experience and remains fully functional even when the user has no internet connectivity.

A second objective is to implement a bidirectional synchronization engine that keeps the local database and the remote cloud database in a consistent state across multiple devices. The synchronization must operate through two channels: a REST-based incremental push and pull mechanism for batch data transfer, and a WebSocket-based real-time channel for character-by-character propagation of live edits to all active sessions. When the same note is modified on two different devices before a sync cycle completes, the system must resolve the conflict deterministically using a last-write-wins strategy based on millisecond-precision timestamps.

A third objective is to architect and implement the Host-Your-Own-Data model, which enables any user to deploy their own independent, private backend on their personal AWS account. The deployment must be achievable through a single AWS CloudFormation template that the user downloads from within the application and uploads to the AWS Console. The provisioned stack must include Amazon Cognito for authentication, API Gateway for REST and WebSocket endpoints, Lambda functions for compute logic, a DynamoDB table for data persistence, and an S3 bucket for asset storage. At no point should the user be required to share their AWS credentials with NoteCognition or any third party.

A fourth objective is to build a rich, modern frontend experience. The application must include a split-pane Markdown editor with live-rendered preview, a hierarchical folder tree with drag-and-drop reorganization, inline image attachments via drag-and-drop or clipboard paste, formatting toolbar shortcuts, and a polished dark-mode glassmorphism user interface with responsive layouts and smooth micro-animations. The user interface must feel premium, performant, and intuitive on both desktop and tablet screen sizes.

Finally, the project aims to validate that the entire system can operate within the AWS Free Tier, making it economically accessible to individual users and students without incurring significant cloud infrastructure expenses.

\section{Purpose}

The purpose of this project is to provide a practical, working solution to a real-world problem that affects anyone who takes notes digitally: the loss of control over personal data when using cloud-hosted platforms. Every day, millions of users type their thoughts, lecture notes, research findings, project plans, and personal reflections into applications that store this information on servers they do not own, in jurisdictions they may not be aware of, under terms of service they rarely read. The implicit bargain, convenience in exchange for data custody, is one that an increasing number of users are no longer comfortable making.

NoteCognition exists to demonstrate that this trade-off is not necessary. A well-engineered application can give users the full convenience of cloud synchronization and real-time cross-device editing while ensuring that their data never touches a server they do not personally own and administer. By building the backend as a portable CloudFormation template and making it deployable on the user's own AWS account, NoteCognition shifts the locus of control from the application developer to the end user. The user decides which AWS region hosts their data, the user holds the encryption keys, the user controls the billing, and the user can delete the entire backend with a single stack deletion command at any time.

Beyond the privacy argument, this project also serves an educational purpose. It demonstrates the practical application of several important concepts from modern software engineering: local-first application design, offline-first progressive web architecture, serverless computing, infrastructure as code, single-table DynamoDB design patterns, WebSocket-based real-time communication, and JWT-based stateless authentication. The project serves as a comprehensive case study in how these technologies can be integrated into a cohesive, production-quality system.

\section{Scope}

The scope of this project encompasses the design, development, testing, and deployment of the complete NoteCognition platform, including both the frontend client application and the serverless AWS backend infrastructure.

On the frontend side, the project covers the implementation of the local IndexedDB database layer using Dexie.js, the Markdown editor with split-pane live preview, the hierarchical folder and file tree with drag-and-drop support, the authentication modal with sign-up, login, and email confirmation flows, the AWS configuration wizard for Host-Your-Own-Data setup, the synchronization engine with offline queue management, the WebSocket client for real-time updates, and the complete user interface with dark-mode glassmorphism styling and responsive layouts. The frontend is built as a single-page application using React 18 and bundled with Vite for optimized development and production builds.

On the backend side, the project covers the design and implementation of the AWS CloudFormation template that defines the entire infrastructure stack, including the Cognito User Pool and App Client for user authentication, the REST API Gateway with Lambda-backed sync endpoints for pulling and pushing data, the WebSocket API Gateway with Lambda-backed connection lifecycle handlers and message broadcast logic, the DynamoDB table with a single-table design optimized for efficient queries, and the S3 bucket for note attachment storage. All Lambda functions are written in Node.js using ES Module syntax and operate under least-privilege IAM roles.

The project also includes the development of PowerShell deployment scripts that automate the packaging of Lambda code, the creation of the bootstrap S3 bucket, the deployment of the CloudFormation stack, and the generation of the local environment configuration file. The project does not include the development of native mobile applications, end-to-end encryption of note contents, or collaborative multi-user editing, although these are identified as future enhancements.

\section{Applicability}

NoteCognition is applicable to a wide range of users and use cases across educational, professional, and personal domains. In the academic setting, students and researchers can use the platform to take lecture notes, organize study materials, draft research papers in Markdown format, and maintain knowledge bases that persist across semesters and devices. The offline-first design is particularly valuable in environments where internet access is unreliable or restricted, such as examination halls, remote field study locations, or campus areas with poor Wi-Fi coverage.

In professional environments, software developers, technical writers, and knowledge workers can use NoteCognition as a lightweight, distraction-free writing tool that supports full Markdown syntax including code blocks, tables, and inline images. The split-pane editor with live preview makes it especially suitable for drafting technical documentation, API reference guides, project specifications, and meeting notes. Because the application runs entirely in the browser and stores data locally, it can be used in corporate environments where data residency and compliance policies prohibit the use of third-party cloud note services.

For privacy-conscious individuals, the Host-Your-Own-Data architecture makes NoteCognition uniquely suitable. Journalists, lawyers, healthcare professionals, and anyone handling sensitive information can deploy the backend on their own AWS account, ensuring that their notes are stored in a region and under a compliance framework of their choosing. The absence of any centralized data collection point means that there is no single server to breach and no third party that can be compelled to hand over user data.

The platform is also applicable as an educational tool for computer science students learning about modern web development, serverless architecture, and infrastructure as code. The complete source code, CloudFormation templates, and deployment scripts serve as a reference implementation of a real-world full-stack application that integrates React, IndexedDB, AWS Lambda, DynamoDB, Cognito, API Gateway, S3, and WebSockets into a cohesive system. Instructors can use the project as a case study for courses on cloud computing, web application development, or distributed systems.

% ===========================
% CHAPTER 2: LITERATURE SURVEY
% ===========================
\chapter{Literature Survey}

\section{Existing Systems}

A comprehensive understanding of the note-taking application landscape requires examining the design philosophies, storage models, and syncing architectures of the most prominent solutions currently available. These platforms represent different trade-offs along the spectrum of user convenience, collaboration capabilities, offline accessibility, and data privacy.

\subsection{Notion}
Notion has emerged as one of the most popular personal and team knowledge management tools, utilizing a rich document model where every paragraph, list item, or image is treated as an independent content block. Notion's primary strength lies in its highly flexible database views and collaboration features. However, its architecture is fundamentally cloud-first and centralized. User documents are stored on Notion's servers (primarily hosted on AWS), meaning that users have zero direct ownership or physical custody of their data. Furthermore, Notion's offline support is notoriously weak; the application requires an active internet connection to load new pages, search notes, or perform any database operations. If the service experiences downtime, users are locked out of their knowledge bases.

\subsection{Obsidian}
Obsidian stands at the opposite end of the spectrum, prioritizing local data ownership. It operates directly on a folder of local Markdown files (a ``vault'') stored on the user's hard drive. Because notes are stored as plain text files, users retain complete ownership, can easily back them up, and can open them with any text editor. Obsidian functions entirely offline with zero latency. However, its main limitation is synchronization across multiple devices. The default solution, Obsidian Sync, is a paid proprietary subscription service. Free alternatives require users to set up complex third-party tools such as Syncthing, Git repositories, or cloud storage clients (e.g., iCloud, Dropbox), which are prone to sync conflicts, high setup friction, and configuration issues. Obsidian also lacks out-of-the-box support for real-time collaboration or live multi-device character synchronization.

\subsection{Evernote}
As one of the pioneers in the digital note-taking space, Evernote popularized the concept of a multi-platform digital notebook. It uses a rich-text formatting model (ENML) and stores all user data on its centralized cloud servers. Over the years, Evernote has faced significant criticism for its bloated client applications, escalating subscription costs, and a major security breach in 2013 that compromised millions of accounts. While it offers basic offline notebook caching on desktop and mobile clients, the underlying data remains fully under the control of the service provider, making it unsuitable for users with strict privacy requirements or data sovereignty concerns.

\subsection{Standard Notes}
Standard Notes is a security-focused application designed specifically to protect user privacy. It achieves this by encrypting all notes client-side before they are transmitted to the sync server, ensuring that not even the service providers can read the user's data. While Standard Notes offers excellent privacy, it suffers from usability constraints. The free version has extremely limited formatting features, and advanced features (such as Markdown editors and folders) are locked behind premium subscriptions. Furthermore, self-hosting a Standard Notes backend is a complex administrative task, requiring the deployment and maintenance of Docker containers, SQL databases, and web servers on a continuous virtual machine, which is costly and time-consuming for non-technical users.

\subsection{Joplin}
Joplin is an open-source note-taking application that supports Markdown and local-first storage. It offers optional end-to-end encryption and supports syncing notes to a variety of cloud providers, including Dropbox, OneDrive, WebDAV, or its own proprietary Joplin Cloud. While Joplin addresses many privacy concerns, it lacks built-in real-time WebSocket-based synchronization. In addition, its self-hosting options (such as Joplin Server) require running a persistent node server and database, incurring constant monthly server maintenance and hosting costs, rather than leveraging serverless pay-per-use resources.

\section{Comparison of Existing Systems}

The table below presents a comparative analysis of these platforms against NoteCognition, illustrating the unique combination of local-first speed, open-source transparency, one-click serverless self-hosting, and real-time synchronization that NoteCognition introduces.

\begin{table}[H]
\centering
\caption{Comparison of Existing Note-Taking Platforms}
\label{tab:comparison}
\begin{tabularx}{\textwidth}{|l|X|X|X|X|X|}
\hline
\textbf{Feature} & \textbf{Notion} & \textbf{Obsidian} & \textbf{Evernote} & \textbf{Joplin} & \textbf{NoteCognition} \\
\hline
Offline Support & Partial & Full & Partial & Full & Full \\
\hline
Open Source & No & No & No & Yes & Yes \\
\hline
Self-Hosted Backend & No & No & No & Partial & Yes (1-Click) \\
\hline
Real-Time Sync & Yes & No & Yes & No & Yes (WebSocket) \\
\hline
Markdown Support & Limited & Full & No & Full & Full \\
\hline
Data Ownership & No & Local Files & No & Partial & Full \\
\hline
\end{tabularx}
\end{table}

\section{Technology Review}

NoteCognition is built on a modern, highly optimized stack that combines web technologies on the frontend with serverless infrastructure on the backend. This section reviews the key technologies utilized in the project.

\subsection{React and Vite}
The client application is built using React 18, a declarative component-based UI library. React's virtual DOM reconciliation ensures high-performance rendering of note lists, search results, and complex folder structures. Vite is used as the frontend build tool and bundler. Unlike traditional bundlers like Webpack, Vite leverages native ES modules in the browser during development for instantaneous hot module replacement (HMR) and uses Rollup for highly optimized production builds.

\subsection{Dexie.js (IndexedDB)}
To achieve local-first performance, NoteCognition uses IndexedDB, a low-level API for client-side storage of large amounts of structured data. Because the native IndexedDB API is event-driven and complex, the project integrates Dexie.js, a lightweight wrapper that provides a developer-friendly, Promise-based API. Dexie.js simplifies query construction, transaction management, and schema versioning, and provides hooks like \texttt{useLiveQuery} to make UI updates reactive to local database modifications.

\subsection{AWS Cognito}
User authentication, token management, and secure registration flows are managed using Amazon Cognito User Pools. Cognito provides a secure, fully managed identity store that handles sign-up verification emails, password hashing, and token issuance (JWT) without requiring the development team to build or maintain sensitive authentication databases.

\subsection{AWS API Gateway (REST and WebSockets)}
Amazon API Gateway is utilized to expose the backend APIs. A REST API Gateway routes batch synchronization requests (pulling and pushing changes) to Lambda compute services. In parallel, a WebSocket API Gateway manages persistent bidirectional TCP connections, allowing the server to broadcast changes to connected client sessions within milliseconds.

\subsection{AWS Lambda}
AWS Lambda is a serverless, event-driven compute service that executes the application's business logic. Written in Node.js, the Lambda functions are stateless, scale automatically with incoming traffic, and incur zero costs when the application is idle.

\subsection{Amazon DynamoDB}
For persistent cloud storage, NoteCognition uses Amazon DynamoDB, a fully managed NoSQL database. The project implements a Single-Table Design pattern, where all entity types (users, folders, notes, and websocket connections) are co-located in a single table. This pattern minimizes the number of roundtrips to the database and allows efficient fetching of related entities in a single query using composite primary keys.

\subsection{Amazon S3}
Amazon Simple Storage Service (S3) provides highly durable object storage for note attachments, such as images. S3 handles large binary payloads efficiently, offloading file storage from the primary DynamoDB database.

\subsection{AWS CloudFormation}
AWS CloudFormation is an Infrastructure as Code (IaC) service that allows the entire serverless backend stack to be defined declaratively in a single YAML template. This template enables users to deploy their private backend with a single click in their personal AWS account.

\section{Gaps in Existing Systems}

An analysis of the existing literature and note-taking systems reveals a critical gap: users are forced to choose between the seamless synchronization and real-time features of centralized cloud platforms (e.g., Notion) and the strict data privacy and local speed of offline-first filesystems (e.g., Obsidian). No popular, open-source note-taking application offers a built-in, low-cost, zero-maintenance syncing backend that can be deployed into a user's private cloud account with a single click.

NoteCognition addresses this gap by combining a local-first browser architecture with a self-deployable serverless AWS backend. By treating IndexedDB as the primary source of truth, it guarantees offline usability and instant responsiveness. By utilizing CloudFormation, it eliminates the technical barriers to self-hosting, permitting non-technical users to own their cloud backend. By leveraging AWS serverless technologies (Lambda, API Gateway, DynamoDB, Cognito, S3), it ensures that the private backend is completely free of ongoing server maintenance and operates within the AWS Free Tier, bridging the gap between convenience, cost, and data sovereignty.

% ===========================
% CHAPTER 3: SYSTEM ANALYSIS
% ===========================
\chapter{System Analysis}

\section{Requirement Analysis}

\subsection{Functional Requirements}
% TODO: Add detailed functional requirements
\begin{enumerate}
    \item The system shall allow users to create, read, update, and delete notes in Markdown format.
    \item The system shall store all data locally using IndexedDB for offline-first access.
    \item The system shall synchronize data in real-time across devices using WebSockets.
    \item The system shall provide user authentication via AWS Cognito.
    \item The system shall allow users to organize notes in a hierarchical folder structure.
    \item The system shall support drag-and-drop image uploads to notes.
    \item The system shall enable one-click deployment of a private backend using CloudFormation.
    \item The system shall implement conflict resolution using a last-write-wins strategy.
\end{enumerate}

\subsection{Non-Functional Requirements}
% TODO: Add non-functional requirements
\begin{enumerate}
    \item \textbf{Performance:} The editor must respond to input within 16ms (60 FPS) for a lag-free experience.
    \item \textbf{Availability:} The application must be fully functional offline with local data persistence.
    \item \textbf{Scalability:} The serverless backend must auto-scale with user demand.
    \item \textbf{Security:} All API endpoints must validate JWT tokens; no AWS credentials stored client-side.
    \item \textbf{Usability:} The interface must follow modern UI/UX design principles with responsive layouts.
\end{enumerate}

\section{Feasibility Study}

A feasibility study is essential to evaluate the practicality, cost, and operational viability of NoteCognition before proceeding with development. The study is divided into technical, economic, and operational dimensions.

\subsection{Technical Feasibility}

The technical feasibility of NoteCognition hinges on the maturity of browser-based databases and serverless cloud computing. On the client side, modern browsers support IndexedDB, a robust, transactional NoSQL database capable of storing gigabytes of structured note and folder data locally with sub-millisecond retrieval times. Libraries such as Dexie.js provide high-level abstractions that make transaction management and database reactivity simple and reliable. On the server side, AWS serverless technologies (Cognito, Lambda, API Gateway, DynamoDB, and S3) provide fully managed services that scale automatically and do not require server maintenance. The dynamic self-hosting aspect is made technically feasible by AWS CloudFormation, which allows infrastructure configurations to be encapsulated into a single downloadable template. The WebSocket protocol is fully supported by modern browsers and API Gateway, making real-time, low-latency cross-device synchronization feasible. Thus, the project is technically feasible.

\subsection{Economic Feasibility}

The economic feasibility of a software project is determined by its development and operational costs. NoteCognition is highly cost-efficient due to its serverless architecture and host-your-own-data model. Since users host their own backend on their personal AWS accounts, the development team does not incur ongoing hosting, compute, or database fees. For individual users, AWS offers a generous Free Tier:
\begin{itemize}[noitemsep]
    \item \textbf{Amazon Cognito:} Up to 50,000 monthly active users free.
    \item \textbf{AWS Lambda:} 1 million free requests and 400,000 GB-seconds of compute time per month.
    \item \textbf{Amazon DynamoDB:} 25 GB of storage and 25 write/read capacity units free.
    \item \textbf{Amazon S3:} 5 GB of standard storage free.
    \item \textbf{Amazon API Gateway:} 1 million REST API calls and 1 million WebSocket connection minutes free.
\end{itemize}
Because NoteCognition's resource consumption for a single user is well within these limits, the operational cost for the end user is virtually zero. On the client side, the frontend is built using free, open-source technologies (React, Vite, Dexie.js) and hosted on Netlify's free starter tier. Therefore, the project is economically feasible.

\subsection{Operational Feasibility}

Operational feasibility assesses how well the system will be accepted by users and integrated into their workflows. NoteCognition addresses the two key issues in digital note-taking: usability and privacy. The frontend provides a familiar, intuitive Markdown editing environment with a live preview pane, standard formatting shortcuts, and a responsive folder structure, minimizing the user learning curve. For privacy-conscious users, the host-your-own-data setup is simplified into a step-by-step wizard. Users do not need cloud administration skills or access keys; they simply upload a template to the AWS Console and copy-paste five output configuration strings. From a maintenance perspective, the serverless backend requires zero patching, upgrading, or capacity planning, ensuring high operational reliability. Therefore, the project is operationally feasible.

\section{Software and Hardware Requirements}

\subsection{Software Requirements}
\begin{table}[H]
\centering
\caption{Software Requirements}
\label{tab:software-req}
\begin{tabular}{|l|l|}
\hline
\textbf{Component} & \textbf{Technology / Tool} \\
\hline
Frontend Framework & React 18.3 with Vite 6.3 \\
\hline
Local Database & Dexie.js 4.4 (IndexedDB) \\
\hline
UI Components & Radix UI, MUI, Lucide Icons \\
\hline
Styling & TailwindCSS 4.1, Vanilla CSS \\
\hline
State Management & React Hooks (useState, useEffect, useLiveQuery) \\
\hline
Authentication SDK & amazon-cognito-identity-js 6.3 \\
\hline
Cloud Provider & Amazon Web Services (AWS) \\
\hline
Backend Compute & AWS Lambda (Node.js 18.x runtime) \\
\hline
Database (Cloud) & Amazon DynamoDB \\
\hline
API Layer & AWS API Gateway (REST + WebSocket) \\
\hline
File Storage & Amazon S3 \\
\hline
Infrastructure as Code & AWS CloudFormation \\
\hline
Version Control & Git / GitHub \\
\hline
Package Manager & npm \\
\hline
IDE & Visual Studio Code \\
\hline
\end{tabular}
\end{table}

\subsection{Hardware Requirements}
\begin{table}[H]
\centering
\caption{Hardware Requirements}
\label{tab:hardware-req}
\begin{tabular}{|l|l|}
\hline
\textbf{Component} & \textbf{Minimum Specification} \\
\hline
Processor & Intel Core i3 / AMD Ryzen 3 or higher \\
\hline
RAM & 4 GB (8 GB recommended) \\
\hline
Storage & 500 MB free disk space \\
\hline
Internet & Required for sync \& cloud deployment \\
\hline
Browser & Chrome 90+, Firefox 88+, Edge 90+, Safari 14+ \\
\hline
\end{tabular}
\end{table}

\section{Project Plan and Timeline}

The development of NoteCognition was structured across a 16-week timeline, divided into seven distinct phases. This structured approach ensured that the core local-first capabilities were fully validated before introducing cloud synchronization and user dynamic configuration features.

\subsection{Timeline Chart}

The table below outlines the timeline and phase-wise progression of the project.

\begin{table}[H]
\centering
\caption{Project Timeline and Phase-wise Allocation}
\label{tab:timeline}
\begin{tabularx}{\textwidth}{|l|p{3.5cm}|X|}
\hline
\textbf{Phase} & \textbf{Timeline} & \textbf{Key Tasks and Deliverables} \\
\hline
Phase 1 & Week 1 -- Week 3 & Requirement Analysis, Literature Survey, System Architecture Definition \\
\hline
Phase 2 & Week 4 -- Week 6 & Frontend Environment Setup, IndexedDB Schema Setup (Dexie.js), Markdown Editor Implementation \\
\hline
Phase 3 & Week 7 -- Week 9 & AWS CloudFormation Template Design (Cognito, Lambda, DynamoDB, S3), PowerShell Deploy Scripts \\
\hline
Phase 4 & Week 10 -- Week 11 & Incremental Sync Logic (REST), WebSocket Real-Time Handlers, Image Upload Service \\
\hline
Phase 5 & Week 12 -- Week 13 & Host-Your-Own-Data Integration, Setup Wizard UI Modal, LocalStorage Refactoring \\
\hline
Phase 6 & Week 14 -- Week 15 & System Integration Testing, Security Audit, Bug Fixes, Sync Conflict Resolution Edge Cases \\
\hline
Phase 7 & Week 16 & Live Deployment (Netlify), Verification against Live AWS Stack, Project Report Submission \\
\hline
\end{tabularx}
\end{table}

\subsection{Phase-wise Implementation Details}

\begin{description}
    \item[Phase 1 — Requirement Analysis and Research (Weeks 1 to 3):] This phase was dedicated to conducting a detailed literature review of existing note-taking tools, identifying their design limitations (centralized servers or complex syncing), and defining the system specifications. The project scope, technical stack, and software/hardware requirements were finalized during this time.
    \item[Phase 2 — Frontend Foundation and Local DB Setup (Weeks 4 to 6):] We initialized the frontend codebase using React 18 and Vite. The primary IndexedDB database and tables were configured using Dexie.js. The user interface layout, including the responsive sidebar folder tree, toolbar shortcuts, status bar, and split-pane Markdown editor with live preview rendering, was fully implemented to operate locally with zero network reliance.
    \item[Phase 3 — AWS Serverless Backend Design (Weeks 7 to 9):] The backend serverless resources were designed and defined. We authored the AWS CloudFormation template declaring Amazon Cognito for user authentication, Amazon API Gateway (REST and WebSocket routes), AWS Lambda functions for backend logic, and DynamoDB for persistent storage. PowerShell deployment scripts were developed to automate packaging and stack deployment.
    \item[Phase 4 — Sync Engine and Real-Time WebSocket (Weeks 10 to 11):] We implemented the bidirectional synchronization engine on the client (SyncService) and server (Sync Lambda). A timestamp-based pull and incremental push mechanism was integrated. Real-time updates were enabled via WebSocket connections, allowing live edits to propagate character-by-character across multiple active browser sessions of the same user.
    \item[Phase 5 — Host-Your-Own-Data Integration (Weeks 12 to 13):] To support private deployments, the frontend services were modified to dynamically read connection settings from browser localStorage rather than static environment variables. We built the AWS Configuration Wizard UI modal to guide users through downloading the CloudFormation template, deploying it on their AWS Console, and configuring the connection.
    \item[Phase 6 — Testing and Bug Fixes (Weeks 14 to 15):] Comprehensive testing was conducted across all modules. Functional tests verified note CRUD, folder trees, authentication, sync, and the host-your-own workflow. Edge cases in the sync conflict resolution logic were tested, and UI responsiveness and animations were polished.
    \item[Phase 7 — Deployment and Documentation (Week 16):] The frontend was deployed to production on Netlify. The deployment was verified against a live AWS backend stack. The project report and documentation were prepared and compiled for final academic submission.
\end{description}

% ===========================
% CHAPTER 4: SYSTEM DESIGN
% ===========================
\chapter{System Design}

\section{System Architecture}

NoteCognition utilizes a Local-First, Serverless Sync architecture designed to separate client-side responsiveness from server-side persistence. The frontend browser layer handles all immediate user interactions locally, while the serverless AWS backend serves as a secure, synchronized replication log.

The primary data store is the browser's IndexedDB database, managed via the Dexie.js library. This local database handles all read and write transactions. Because every user interaction is executed directly against this local database, the editor achieves zero perceptible latency. If the client is offline, modifications are written to IndexedDB and appended to a persistent synchronization queue. When the client detects that network connectivity has been restored, it flushes the queue, transmitting changes to the backend in batches.

The backend infrastructure is built entirely using serverless AWS resources. An API Gateway exposes REST and WebSocket endpoints. The REST endpoints handle authentication, initial database synchronization, and conflict resolution by invoking AWS Lambda functions. The WebSocket endpoints manage persistent connections to propagate character-by-character updates to all logged-in devices of the user. Amazon DynamoDB persistent storage holds notes, folders, and connection metadata. Amazon S3 holds media assets (images, attachments). Amazon Cognito handles registration, logins, and token validation.

\begin{figure}[H]
    \centering
    \includegraphics[width=0.9\textwidth]{images/system_architecture.png}
    \caption{NoteCognition Local-First Serverless System Architecture}
    \label{fig:system_architecture}
\end{figure}

\section{Methodology and Algorithms}

To maintain data integrity and consistency across multiple offline-capable devices, NoteCognition implements specific client-server synchronization and conflict resolution algorithms.

\subsection{Incremental Synchronization Algorithm}

The synchronization process follows an incremental sync loop:
\begin{enumerate}[noitemsep]
    \item \textbf{Client Pull Phase:} Upon establishing a session or recovering internet connection, the client sends a GET request to the REST API `/sync/pullAll` endpoint with its `lastSyncTimestamp`. The Lambda function queries DynamoDB for all note and folder changes modified by this user after the client's `lastSyncTimestamp`. The client merges these updates into its local IndexedDB database.
    \item \textbf{Client Push Phase:} If there are modifications pending in the client's sync queue, the client pushes the batch of modifications as a JSON payload to the REST API `/sync/push` endpoint. The Lambda function performs bulk updates in DynamoDB.
    \item \textbf{Timestamp Synchronization:} Once the pull and push phases complete successfully, the client updates its local `lastSyncTimestamp` to the current server-returned timestamp, preparing for the next incremental cycle.
\end{enumerate}

\subsection{Last-Write-Wins (LWW) Conflict Resolution}

When a user modifies the same note on two separate devices while offline, a sync conflict occurs when both devices reconnect. NoteCognition resolves this conflict deterministically using a Last-Write-Wins (LWW) conflict resolution strategy:
\begin{itemize}[noitemsep]
    \item Every database record (Note or Folder) contains an `updatedAt` field tracking the epoch time of its last modification in milliseconds.
    \item When merging remote changes into the local IndexedDB or applying incoming modifications to DynamoDB, the system compares the incoming record's `updatedAt` timestamp against the existing record's `updatedAt` timestamp.
    \item If the incoming record's timestamp is strictly greater than the local record's timestamp, the incoming record overwrites the local record. Otherwise, the incoming record is discarded.
\end{itemize}

\section{Data Flow Diagram (DFD)}

Data Flow Diagrams (DFDs) trace the flow of data through NoteCognition's logical components and data stores.

\subsection{Context Diagram (Level 0 DFD)}

The Context Diagram defines the system boundary of NoteCognition. The system interacts with three external entities: the User, Amazon Cognito, and the AWS Cloud Services. The User provides input (notes, folder structures, credentials) and receives visual feedback. Cognito manages user registration and authentication, and the AWS Cloud Services handle persistent storage and cross-device sync.

\begin{figure}[H]
    \centering
    \includegraphics[width=0.8\textwidth]{images/dfd_level_0.png}
    \caption{NoteCognition Context Diagram (Level 0 DFD)}
    \label{fig:dfd_level_0}
\end{figure}

\subsection{Level 1 DFD}

The Level 1 DFD decomposes the system boundary into five core processes:
\begin{enumerate}[noitemsep]
    \item \textbf{Note Management Process:} Reads and writes note data to the local IndexedDB data store.
    \item \textbf{Folder Management Process:} Manages the hierarchical organization of folders in IndexedDB.
    \item \textbf{User Authentication Process:} Integrates with Cognito to manage user registration, login, and JWT verification.
    \item \textbf{Sync Engine Process:} Evaluates the offline queue and coordinates REST pull/push calls to sync with DynamoDB.
    \item \textbf{Real-Time WebSocket Process:} Listens for active connections and broadcasts delta edits to sync sessions in real-time.
\end{enumerate}

\begin{figure}[H]
    \centering
    \includegraphics[width=0.9\textwidth]{images/dfd_level_1.png}
    \caption{NoteCognition Level 1 Data Flow Diagram}
    \label{fig:dfd_level_1}
\end{figure}

\section{Flowchart Diagram}

The flowchart traces the complete lifecycle of NoteCognition, starting from application launch, through user authentication, local database initialization, WebSocket connection setup, the editor rendering loop, offline queue tracking, and cloud replication.

\begin{figure}[H]
    \centering
    \includegraphics[width=0.8\textwidth]{images/flowchart.png}
    \caption{NoteCognition Application Lifecycle Flowchart}
    \label{fig:flowchart}
\end{figure}

\section{Use Case Diagram}

The Use Case Diagram defines the capabilities available to the User and their relationships to the underlying system services. The User interacts with note use cases (Create Note, Edit Note, Delete Note, Export Markdown), folder use cases (Manage Folders, Drag-Drop Reorder), authentication use cases (Sign Up, Log In), and self-hosting use cases (Upload Image Attachment, Deploy Private Backend, Configure AWS settings). The backend services (Cognito, DynamoDB, S3, CloudFormation) act as supporting actors that fulfill these use cases.

\begin{figure}[H]
    \centering
    \includegraphics[width=0.9\textwidth]{images/use_case_diagram.png}
    \caption{NoteCognition Use Case Diagram}
    \label{fig:use_case}
\end{figure}

\section{Sequence Diagrams}

Sequence diagrams illustrate the chronological flow of messages exchanged between system components for critical operations.

\subsection{User Authentication Flow}

When a user registers or logs in, the client communicates directly with the Cognito User Pool endpoint. Upon successful verification, Cognito issues a set of JSON Web Tokens (JWT) including an ID token, Access token, and Refresh token. These tokens are cached locally and attached to subsequent API Gateway request headers for stateless authorization on the backend.

\subsection{Note Synchronization and Real-Time WebSocket Propagation}

When a user edits a note on Device A, the change is written locally to IndexedDB. If online, the client pushes the modification to API Gateway, which invokes the Sync Lambda to update DynamoDB. Simultaneously, a WebSocket message is sent through API Gateway WebSockets, which reads active connection records from DynamoDB and broadcasts the edit to all other registered connections belonging to that user. Device B receives the WebSocket broadcast, merges the changes using the LWW strategy, and dynamically updates its UI.

\begin{figure}[H]
    \centering
    \includegraphics[width=0.9\textwidth]{images/sequence_diagram.png}
    \caption{Note Modification and Real-Time Synchronization Sequence}
    \label{fig:sequence_diagram}
\end{figure}

\subsection{Host-Your-Own-Data Stack Deployment Flow}

The host-your-own-data workflow decouples the deployment process from backend code. The user clicks ``Cloud Host'' inside the application to fetch the CloudFormation template (`notecognition-template.yaml`) served from S3/Netlify. The client redirects the user to the AWS Console in a new tab, where the user deploys the template. The stack creates the resources in the user's personal account and outputs connection strings (`ApiUrl`, `UserPoolId`, etc.). The user copy-pastes these values into the client UI, which stores them in `localStorage` and routes all API operations to the user's private cloud endpoints.

\section{Activity Diagram}

The Activity Diagram captures the operational workflow of the Synchronization Service. It evaluates connection availability, manages pulls, resolves conflicts, processes queue flushes, handles API communication errors with retries, and clears local logs upon success.

\begin{figure}[H]
    \centering
    \includegraphics[width=0.7\textwidth]{images/activity_diagram.png}
    \caption{NoteCognition Client-Server Sync Activity Diagram}
    \label{fig:activity_diagram}
\end{figure}

\section{Class Diagram}

The Class Diagram illustrates the object-oriented structure of the NoteCognition frontend application. It displays core models (`Note`, `Folder`), database configurations (`DexieDB`), services (`AuthService`, `SyncService`, `ConfigService`, `UploadService`), and React controller components (`EditorArea`, `FileTree`, `App`).

\begin{figure}[H]
    \centering
    \includegraphics[width=0.9\textwidth]{images/class_diagram.png}
    \caption{NoteCognition Frontend Class Diagram}
    \label{fig:class_diagram}
\end{figure}

\section{Database Design}

\subsection{Local Database Schema (IndexedDB)}
The local IndexedDB managed via Dexie.js uses the following schema:

\begin{table}[H]
\centering
\caption{IndexedDB Schema — Notes Table}
\label{tab:indexeddb-notes}
\begin{tabular}{|l|l|l|}
\hline
\textbf{Field} & \textbf{Type} & \textbf{Description} \\
\hline
id & String (UUID) & Primary key \\
\hline
title & String & Note title \\
\hline
body & String & Markdown content \\
\hline
folderId & String & Parent folder reference \\
\hline
isDeleted & Boolean & Soft-delete flag \\
\hline
createdAt & Number & Creation timestamp \\
\hline
updatedAt & Number & Last modification timestamp \\
\hline
\end{tabular}
\end{table}

\subsection{Cloud Database Schema (DynamoDB Single-Table)}
All entities are stored in a single DynamoDB table (\texttt{notecognition-data-dev}) using composite keys:

\begin{table}[H]
\centering
\caption{DynamoDB Single-Table Schema}
\label{tab:dynamodb-schema}
\begin{tabularx}{\textwidth}{|l|l|l|X|}
\hline
\textbf{Entity} & \textbf{PK} & \textbf{SK} & \textbf{Attributes} \\
\hline
User Account & \texttt{USER\#<Id>} & \texttt{METADATA} & Email, Signup Timestamp \\
\hline
Folder & \texttt{USER\#<Id>} & \texttt{FOLDER\#<FId>} & Name, ParentFolderId, UpdatedAt \\
\hline
Note & \texttt{USER\#<Id>} & \texttt{NOTE\#<NId>} & Title, Body, FolderId, UpdatedAt, isDeleted \\
\hline
Socket & \texttt{USER\#<Id>} & \texttt{CONN\#<CId>} & ConnectionId, ConnectedAt \\
\hline
\end{tabularx}
\end{table}

\section{Component Design}

The physical components of NoteCognition are organized into a tiered architecture that divides presentation, business logic, client data, and external services:
\begin{itemize}[noitemsep]
    \item \textbf{UI Layer (React):} Components like `App`, `Sidebar`, `FileTree`, `EditorArea`, and dialog modals manage rendering and handle user events.
    \item \textbf{Service Layer:} Modules like `AuthService` (Cognito integration), `SyncService` (REST and WebSocket synchronization), `UploadService` (S3 attachments), and `ConfigService` (localStorage AWS configuration) encapsulate business logic.
    \item \textbf{Local Storage Layer:} `IndexedDB` and `localStorage` manage cached client-side state.
    \item \textbf{Cloud Services Layer:} Cognito, API Gateway (REST and WS), AWS Lambda, DynamoDB, and S3 process and persist sync data.
\end{itemize}

\begin{figure}[H]
    \centering
    \includegraphics[width=0.9\textwidth]{images/component_diagram.png}
    \caption{NoteCognition Logical Component Diagram}
    \label{fig:component_diagram}
\end{figure}

\section{Deployment Diagram}

The Deployment Diagram defines the physical deployment nodes of NoteCognition:
\begin{enumerate}[noitemsep]
    \item \textbf{Client Node (Browser):} Executes the compiled React SPA code and runs the local IndexedDB database.
    \item \textbf{Netlify CDN Node:} Hosts and serves the static HTML, JavaScript, CSS, and CloudFormation template assets.
    \item \textbf{AWS Cloud Account Node:} Encompasses Cognito for authentication, API Gateway endpoints, Lambda compute functions, DynamoDB persistent table, and the S3 assets bucket.
\end{enumerate}

\begin{figure}[H]
    \centering
    \includegraphics[width=0.9\textwidth]{images/deployment_diagram.png}
    \caption{NoteCognition Deployment Diagram}
    \label{fig:deployment_diagram}
\end{figure}

\section{User Interface Design}

The user interface of NoteCognition is designed around modern UI/UX principles, focusing on high visual appeal, responsive layout grids, and visual clarity. Key design features include:
\begin{itemize}[noitemsep]
    \item \textbf{Glassmorphism Aesthetics:} Leverages semi-transparent backdrop filters, subtle borders, and harmonious dark HSL color palettes to create a modern visual look.
    \item \textbf{Split-Pane Editor:} Features a side-by-side editing layout with a Markdown code area on the left and a live-rendered HTML preview on the right.
    \item \textbf{Hierarchical Navigation:} Provides a nested folder and note directory tree in the sidebar with drag-and-drop capabilities.
    \item \textbf{Connection Status Feedback:} Shows a status indicator (Online/Offline/Syncing) in the bottom status bar, giving users real-time visibility into the sync engine state.
\end{itemize}

% ===========================
% CHAPTER 5: RESULTS AND EXPLANATION
% ===========================
\chapter{Results and Explanation}

\section{Implementation Approaches}
The NoteCognition application is engineered around four core implementation paradigms designed to guarantee offline operation, responsive user interactions, data sovereignty, and efficient synchronization:
\begin{enumerate}
    \item \textbf{Local-First Data Architecture:} Traditional applications block user interactions while waiting for a server request to complete. NoteCognition writes all modifications (note creation, content edits, folder restructuring) directly to the client browser's local database (IndexedDB) via the Dexie.js wrapper. IndexedDB operations complete in less than 5 milliseconds, meaning the user interface experiences zero lag regardless of network connectivity.
    \item \textbf{Serverless Cloud-Native Backend:} The cloud-based backend uses a serverless model. AWS Lambda functions run in response to specific events (HTTP requests from API Gateway or DynamoDB stream events), eliminating the idle costs and maintenance overhead of 24/7 virtual servers. Amazon DynamoDB provides a single-table NoSQL structure that scales partition throughput dynamically.
    \item \textbf{Host-Your-Own-Data (HYOD) Paradigm:} To address cloud data privacy concerns, the client interface provides a configuration modal where users can supply their own AWS credentials and endpoint URLs. When custom values are provided, the application overrides its default service values, stores the credentials securely in \texttt{localStorage}, and directs all API sync calls, WebSocket events, and S3 asset uploads directly to the user's personal AWS account.
    \item \textbf{Incremental Log-Structured Synchronization:} Synchronization is timestamp-driven and differential. Instead of uploading the entire note corpus during a sync cycle, the client queries the backend for changes that have occurred since the last recorded sync time (using the \texttt{since} parameter). Active connections are tracked, and real-time broadcasts are pushed via API Gateway WebSockets only to devices associated with the authenticated user ID.
\end{enumerate}

To establish the local development environment, the toolchain utilizes Node.js (v20.11.0+) and Vite with TypeScript. Dependencies are installed using \texttt{npm install}, resolving packages such as Dexie.js for IndexedDB database wrapping, Radix UI for interface components, TailwindCSS for styling, and Cognito client SDKs for secure remote password authentication. Cloud infrastructure resources are deployed using the AWS Serverless Application Model (SAM) command-line interface.

\section{Pseudo Code}
The client-side event hooks and conflict resolution synchronizer operate based on the following structured algorithms:

\begin{lstlisting}[language=Java, caption={Local Write Hook and Sync Queue Trigger}]
// Triggered on local IndexedDB modification
procedure ON_LOCAL_CHANGE(entity_type, operation, entity_id, data)
    data.updatedAt = CURRENT_TIMESTAMP()
    if operation == "creating" then
        data.version = 1
        IndexedDB.write(entity_type, entity_id, data)
    else if operation == "updating" then
        data.version = (data.version || 1) + 1
        IndexedDB.write(entity_type, entity_id, data)
    else if operation == "deleting" then
        IndexedDB.delete(entity_type, entity_id)
    end if
    
    // Asynchronously queue update task
    schedule_async_task(PUSH_TO_CLOUD, entity_type, operation, entity_id, data)
end procedure
\end{lstlisting}

\begin{lstlisting}[language=Java, caption={Bidirectional Synchronization and LWW Conflict Resolution}]
// Main synchronization and conflict resolution loop
procedure SYNCHRONIZE_DATA()
    if NOT IS_NETWORK_AVAILABLE() then return
    
    token = Auth.GET_BEARER_TOKEN()
    last_pull = LocalStorage.GET("notecognition:lastPull")
    
    // Step 1: Pull remote updates from Cloud
    pull_url = "/sync/pull" + (last_pull ? "?since=" + last_pull : "")
    remote_items = HTTP_GET(pull_url, headers={"Authorization": token})
    
    for each item in remote_items do
        local_item = IndexedDB.GET(item.type, item.id)
        
        if local_item == null then
            // Direct write remote item to local database
            APPLY_REMOTE_WRITE(item)
        else
            // Resolve conflict using Last-Write-Wins (LWW)
            local_time = parse_iso_date(local_item.updatedAt)
            remote_time = parse_iso_date(item.updatedAt)
            
            if remote_time > local_time then
                // Remote update is newer; apply locally
                APPLY_REMOTE_WRITE(item)
            else if local_time > remote_time then
                // Local copy is newer; queue a push to update server
                item.version = max(item.version, local_item.version) + 1
                schedule_async_task(PUSH_TO_CLOUD, item.type, "updating", item.id, local_item)
            else
                // Dynamic version tie-breaker if timestamps are equal
                if item.version > local_item.version then
                    APPLY_REMOTE_WRITE(item)
                end if
            end if
        end if
    end for
    
    // Step 2: Flush pending local queue modifications to Server
    for each pending_task in GET_PENDING_QUEUE() do
        response = HTTP_PUT("/file/" + pending_task.id, body=pending_task.data, headers={"Authorization": token})
        
        if response.status == 409 then
            // Conflict returned by API due to version mismatch
            server_item = response.body.serverItem
            RESOLVE_CONFLICT_IMMEDIATELY(pending_task.data, server_item)
        else if response.status == 200 then
            REMOVE_FROM_QUEUE(pending_task.id)
        end if
    end for
    
    LocalStorage.SET("notecognition:lastPull", CURRENT_TIMESTAMP())
end procedure
\end{lstlisting}

\section{Testing}
The verification process for NoteCognition follows a structured testing strategy to ensure reliability across both offline and online states. Testing note-taking software that features offline capabilities requires validating data consistency, sync queue persistence, and conflict resolution alongside basic UI behaviors. The testing methodology includes:
\begin{enumerate}
    \item \textbf{Unit Testing (Vitest):} Individual utility functions and database schemas are verified inside node execution threads. Upgrading the schema from Version 1 through Version 3 is tested to verify it correctly preserves data and applies default values like \texttt{parentId='ROOT'}.
    \item \textbf{Integration Testing (AWS API Client):} The interaction between the client application and AWS services is validated, confirming that Cognito logins SRP processes operate smoothly, S3 pre-signed URLs allow media storage, and DynamoDB Streams successfully trigger the stream-processor Lambdas.
    \item \textbf{System Testing (Simulated Network Offline/Online):} Complete end-to-end user flows are evaluated on simulated offline environments using Developer Tools. The tests verify that notes created while offline are queued, local updates resolve under 5ms, and changes sync to other connected devices once a WebSocket reconnects.
\end{enumerate}

The test cases executed to verify NoteCognition's capabilities are outlined in Table~\ref{tab:test-cases}.

\begin{longtable}{|p{0.8cm}|p{3cm}|p{4cm}|p{3cm}|p{2cm}|}
\caption{Test Cases for NoteCognition} \label{tab:test-cases} \\
\hline
\textbf{ID} & \textbf{Test Case} & \textbf{Input} & \textbf{Expected Output} & \textbf{Status} \\
\hline
\endfirsthead
\hline
\textbf{ID} & \textbf{Test Case} & \textbf{Input} & \textbf{Expected Output} & \textbf{Status} \\
\hline
\endhead
TC01 & User Registration & Valid email \& password & Account created, confirmation email sent & Pass \\
\hline
TC02 & User Login & Valid credentials & JWT token issued, data synced & Pass \\
\hline
TC03 & Create Note & Click ``New File'' & Empty note created in IndexedDB & Pass \\
\hline
TC04 & Edit Note & Type markdown text & Live preview updates in real-time & Pass \\
\hline
TC05 & Offline Edit & Edit note without internet & Changes saved locally & Pass \\
\hline
TC06 & Online Sync & Reconnect to internet & Queued changes pushed to server & Pass \\
\hline
TC07 & Real-Time Sync & Edit on Device A & Changes appear on Device B via WebSocket & Pass \\
\hline
TC08 & Folder Operations & Create/rename/delete folder & Folder tree updates correctly & Pass \\
\hline
TC09 & Image Upload & Drag-and-drop image & Image uploaded to S3, URL inserted & Pass \\
\hline
TC10 & Host-Your-Own & Deploy CloudFormation stack & Private backend operational & Pass \\
\hline
\end{longtable}

\section{Analysis (graphs/chart)}
To evaluate system performance, metrics were collected and analyzed across different network conditions and database sizes. Table~\ref{tab:perf-metrics} summarizes the system response times for database, network, and parser operations.

\begin{table}[H]
\centering
\caption{System Performance Metrics}
\label{tab:perf-metrics}
\begin{tabular}{lllr}
\toprule
\textbf{Category} & \textbf{Operation} & \textbf{Network Context} & \textbf{Avg Latency (ms)} \\
\midrule
Local Write & Create Note (IndexedDB) & Localhost (Offline) & 1.8 \\
Local Write & Update Note (IndexedDB) & Localhost (Offline) & 2.4 \\
Local Rendering & Markdoc Parsing (10KB note) & Localhost & 4.2 \\
Authentication & Cognito SRP Login & WiFi & 142.0 \\
Sync Path & S3 Pre-signed URL generation & WiFi & 64.0 \\
Sync Path & Fetch Cloud Pull (\texttt{/sync/pull}) & WiFi (No updates) & 82.0 \\
WebSocket Sync & Real-Time Note Update & Broadband (Device A to B) & 210.0 \\
\bottomrule
\end{tabular}
\end{table}

Sync latency and reliability were tested across varying network bandwidths, as shown in Table~\ref{tab:network-latencies}.

\begin{table}[H]
\centering
\caption{Sync Latency under Varying Network Environments}
\label{tab:network-latencies}
\begin{tabular}{lrrr}
\toprule
\textbf{Network Profile} & \textbf{Bandwidth (Down/Up)} & \textbf{Ping RTT (ms)} & \textbf{Avg Sync Time (ms)} \\
\midrule
Broadband WiFi & 100 Mbps / 40 Mbps & 12 & 180 \\
Cellular 4G & 15 Mbps / 5 Mbps & 45 & 380 \\
Cellular 3G & 1.5 Mbps / 0.5 Mbps & 180 & 1,120 \\
Offline Mode & 0 Kbps & N/A & Queued (Offline) \\
\bottomrule
\end{tabular}
\end{table}

The analysis indicates that the local-first architecture keeps the client interface responsive during network degradation:
\begin{itemize}
    \item \textbf{UI Response Time:} Local UI operations (typing, file creation) always resolve in under 5ms, independent of network speed. This compares favorably with traditional cloud-based editors, where poor network conditions can cause noticeable interface lag (100ms to 2000ms).
    \item \textbf{Payload Scale Curve:} Network latency grows linearly with note payload size. Average sync times for text notes remain under 400ms for typical note sizes, while larger attachments (such as embedded images) scale based on available network bandwidth.
    \item \textbf{Queue Convergence:} When transitioning from offline to online, the synchronization engine handles concurrent queued changes in parallel, successfully reconciling state changes within 1.5 seconds.
\end{itemize}

Application screenshots illustrating the Markdown Editor interface, Sidebar folder tree layout, and the Host-Your-Own-Data dynamic credentials setup wizard are shown in Figures~\ref{fig:editor}, \ref{fig:sidebar}, and \ref{fig:modal} respectively.

\begin{figure}[H]
    \centering
    \fbox{\begin{minipage}{0.85\textwidth}\centering\vspace{2cm} [Screenshot: Markdown Editor View - Split Pane] \vspace{2cm}\end{minipage}}
    \caption{NoteCognition Markdown Editor — Split Pane View}
    \label{fig:editor}
\end{figure}

\begin{figure}[H]
    \centering
    \fbox{\begin{minipage}{0.85\textwidth}\centering\vspace{2cm} [Screenshot: Folder Tree Sidebar Navigation] \vspace{2cm}\end{minipage}}
    \caption{NoteCognition Sidebar Navigation and Folder Tree View}
    \label{fig:sidebar}
\end{figure}

\begin{figure}[H]
    \centering
    \fbox{\begin{minipage}{0.85\textwidth}\centering\vspace{2cm} [Screenshot: AWS Configurations and Credentials Modal] \vspace{2cm}\end{minipage}}
    \caption{NoteCognition Host-Your-Own-Data (HYOD) Configuration Modal}
    \label{fig:modal}
\end{figure}

% ===========================
% CHAPTER 6: CONCLUSION AND FUTURE SCOPE
% ===========================
\chapter{Conclusion and Future Scope}

\section{Conclusion}
NoteCognition successfully demonstrates a local-first, privacy-centric approach to note-taking. By combining offline-first client architecture with serverless cloud synchronization, the platform achieves zero-lag performance while maintaining data sovereignty through its Host-Your-Own-Data model.

\section{Future Scope}
\begin{enumerate}
    \item \textbf{End-to-End Encryption:} Encrypt note contents before cloud storage for enhanced security.
    \item \textbf{Collaborative Editing:} Implement CRDT-based conflict-free collaborative editing for shared notes.
    \item \textbf{Mobile Application:} Develop native iOS and Android applications using React Native.
    \item \textbf{Plugin System:} Create an extensible plugin architecture for custom Markdown extensions.
    \item \textbf{AI-Powered Features:} Integrate AI-based note summarization, tagging, and search capabilities.
    \item \textbf{Multi-Cloud Support:} Extend Host-Your-Own-Data to support Azure and Google Cloud Platform.
    \item \textbf{Version History:} Implement note versioning with diff visualization and rollback capability.
\end{enumerate}

% ---------------------------------------------------------------------------
% BACK MATTER
% ---------------------------------------------------------------------------

% ===========================
% REFERENCES
% ===========================
\chapter*{References}
\addcontentsline{toc}{chapter}{References}

\begin{enumerate}[label={[\arabic*]}]
    \item React Documentation, ``Getting Started with React,'' Meta Platforms, Inc. [Online]. Available: \url{https://react.dev/}
    \item Vite Documentation, ``Next Generation Frontend Tooling,'' Evan You et al. [Online]. Available: \url{https://vitejs.dev/}
    \item Dexie.js Documentation, ``A Minimalistic Wrapper for IndexedDB,'' David Fahlander. [Online]. Available: \url{https://dexie.org/}
    \item AWS Lambda Documentation, ``Serverless Compute,'' Amazon Web Services, Inc. [Online]. Available: \url{https://docs.aws.amazon.com/lambda/}
    \item Amazon DynamoDB Developer Guide, ``NoSQL Database Service,'' Amazon Web Services, Inc. [Online]. Available: \url{https://docs.aws.amazon.com/dynamodb/}
    \item Amazon Cognito Developer Guide, ``User Authentication and Authorization,'' Amazon Web Services, Inc. [Online]. Available: \url{https://docs.aws.amazon.com/cognito/}
    \item AWS CloudFormation User Guide, ``Infrastructure as Code,'' Amazon Web Services, Inc. [Online]. Available: \url{https://docs.aws.amazon.com/cloudformation/}
    \item Amazon API Gateway Developer Guide, ``REST and WebSocket APIs,'' Amazon Web Services, Inc. [Online]. Available: \url{https://docs.aws.amazon.com/apigateway/}
    \item M. Kleppmann and A. R. Beresford, ``A Conflict-Free Replicated JSON Datatype,'' \textit{IEEE Transactions on Parallel and Distributed Systems}, vol. 28, no. 10, pp. 2733--2746, 2017.
    \item A. Wiggins, ``The Twelve-Factor App,'' 2017. [Online]. Available: \url{https://12factor.net/}
\end{enumerate}

% ===========================
% APPENDICES
% ===========================
\appendix

\chapter{CloudFormation Template}
The following is the complete AWS CloudFormation/SAM Infrastructure-as-Code template used to deploy NoteCognition's serverless backend resources.
\lstinputlisting[language={}, caption={notecognition-template.yaml}]{../cloudformation/template.yaml}

\chapter{Source Code Listings}
This appendix presents key implementation components of the NoteCognition application, specifically the client-side database schemas and synchronization engines.

\section{Database Schema (db.ts)}
The following TypeScript source listing outlines the IndexedDB schemas and version migration stages implemented via Dexie.js.
\lstinputlisting[language=Java, caption={IndexedDB Schema Definition (src/app/db.ts)}]{../src/app/db.ts}

\section{Sync Service (SyncService.ts)}
The TypeScript source listing below illustrates the bidirectional synchronization client, local-first database hooks, conflict resolution handlers, and connection event lifecycles.
\lstinputlisting[language=Java, caption={Bidirectional Synchronization Service (src/app/services/SyncService.ts)}]{../src/app/services/SyncService.ts}

% ===========================
\end{document}
